% Source: https://artofproblemsolving.com/community/c3h2429919_palindrome
% Stable identity: aops-topic-2429919-q01
A number is divisible by $3$ when its digit sum is.
\begin{itemize}
\item One digit: $3,6,9$, so $3$ numbers.
\item Two digits $\overline{aa}$: digit sum $2a$, so $3\mid a$: $33,66,99$, so $3$ numbers.
\item Three digits $\overline{aba}$: digit sum $2a+b$. For each $a\in\{1,\dots,9\}$, exactly $3$ of $b\in\{1,\dots,9\}$ make $2a+b$ a multiple of $3$: $9\cdot3=27$ numbers.
\item Four digits $\overline{abba}<2021$ with nonzero digits: $a=1$ (with $a=2$ the only one below $2021$ is $2002$, which has zeros). The digit sum $2+2b$ is a multiple of $3$ for $b=2,5,8$: $1221,1551,1881$, so $3$ numbers.
\end{itemize}
In total $3+3+27+3=\boxed{36}$.
